\label{sec:050-signatures}

\subsection{SSH Client Features}
\label{subsec:sshclientfeatures}

\autoref{tab:clients} summarizes the basic features of the \nrClients SSH clients we investigated. The most commonly used crypto libraries were OpenSSL~(8), \putty~(4), and LibreSSL~(2)---a fork of OpenSSL. All other libraries appeared only once. All clients support using SSH agents for client authentication; nine are bundled with an agent implementation. All but three libraries support agent forwarding, and two have activated it by default. With regard to nonce generation, we verified that three clients use deterministic nonces for DSA and ECDSA. For three clients, we could not determine whether deterministic nonces are used, as they do not attempt to authenticate multiple times with the same key in the same connection. However, we can rule out either one using $k_\text{proto}$ or RFC~6979 for nonce generation. Regarding biased nonces, we find that the upper nine bits of nonces in \putty~\puttyvulnversion ECDSA NIST~\pfivetwoone signatures are zero, indicating a 9-bit bias. Apart from \putty, we did not find any nonce bias.

For SSH agents (\autoref{tab:agents}), the most commonly used libraries are OpenSSL~(4), \putty~(3), and Libgcrypt~(2). Three use deterministic nonces, and 1Password does not support (EC)DSA. Some Git clients are bundled with SSH clients. This might indicate a slight bias in our dataset towards keys generated by bundled clients.

\begin{table*}
    \centering
    \caption{Overview of SSH clients with their respective versions and operating systems for evaluation. For each client, the table lists whether source code is publicly available (\faCodeBranch), the client ships with an SSH agent implementation (\faUserSecret) or supports the use of agents for authentication (\faSignIn*), and whether agent forwarding is supported (\faShare). A bolt (\faBolt) indicates that agent forwarding is enabled by default, which we consider dangerous. Additionally, for each implementation, the table includes the determined cryptographic library and nonce scheme. For some clients, capturing two signatures for the same plaintext failed due to disconnects after the first authentication attempt. In such cases, we can only rule out known deterministic schemes.}
    \label{tab:clients}
    \begin{tabular}{llcccccll}
         \toprule
         ~ & ~ & ~ & ~ & \multicolumn{3}{c}{SSH Agent} & ~ & ~ \\ \cmidrule{5-7}
         Name & Version & OS & \faCodeBranch & \faUserSecret & \faSignIn* & \faShare & Cryptographic Library & Nonce Scheme \\
         \midrule
         AbsoluteTelnet & 12.16 & \faWindows & \unsupported & \unsupported & \supported & \faBolt & Crypto++ & \textit{Not} RFC 6979 / $k_\text{proto}$ \\
         AsyncSSH & 2.18.0 & \faLinux & \supported & \unsupported & \supported & \supported & OpenSSL 3.3.2 & Random \\
         Bitvise & 9.42 & \faWindows & \unsupported & \unsupported & \supported & \supported & Windows CNG & \textit{Not} RFC 6979 / $k_\text{proto}$ \\
         Cyberduck & 9.0.1 & \faWindows & \supported & \unsupported & \supported & \unsupported & BC Java 1.78.1 & Random \\
         Dropbear & 2024.86 & \faLinux & \supported & \unsupported & \supported & \supported & LibTomCrypt 1.18.2 & Random \\
         Erlang/OTP SSH & 5.2.1 & \faLinux & \supported & \unsupported & \supported & \unsupported & OpenSSL 3.3.2 & \textit{Not} RFC 6979 / $k_\text{proto}$ \\
         FileZilla & 3.67.0 & \faWindows & \supported & \unsupported & \supported & \unsupported & \putty 0.73 & Random \\
         Golang x/crypto/ssh & 0.29.0 & \faWindows & \supported  & \unsupported & \supported & \supported & Go crypto 1.23.3 & Optional (Any) \\
         libssh & 0.11.1 & \faLinux & \supported & \unsupported & \supported & \supported & OpenSSL 3.3.2 & Random \\
         ~ & ~ & ~ & ~ & ~ & ~ & ~ & MBedTLS 3.6.2 & RFC 6979 \\
         OpenSSH Portable & 9.9p1 & \faLinux & \supported & \supported & \supported & \supported & OpenSSL 3.3.2 & Random \\
         ~ & ~ & \faApple & \supported & \supported & \supported & \supported & LibreSSL 3.3.6 & Random \\
         Paramiko & 3.5.0 & \faLinux & \supported & \unsupported & \supported & \supported & OpenSSL 3.3.2 & Random \\
         PKIX-SSH & 15.3 & \faLinux & \supported & \supported & \supported & \supported & OpenSSL 3.3.2 & Random \\
         PortX & 2.2.12 & \faApple & \unsupported & \unsupported & \supported & \supported & OpenSSL 1.1.1 & Random \\
         \putty & \puttyvulnversion & \faWindows & \supported & \supported & \supported & \supported & \putty \puttyvulnversion & $k_\text{proto}$ \\
         ~ & \puttyfixedversion & \faWindows & \supported & \supported & \supported & \supported & \putty \puttyfixedversion & RFC 6979 \\
         SecureCRT & 9.5.2 & \faWindows & \unsupported & \supported & \supported & \supported & BSAFE Crypto-C ME 4.1.5 & Random \\
         Secure ShellFish & 2025.17 & \faApple & \unsupported & \supported & \supported & \supported & OpenSSL & \textit{Not} RFC 6979 / $k_\text{proto}$ \\
         ServerCat & 1.18 & \faApple & \unsupported & \unsupported & \unsupported & \unsupported & BoringSSL & \textit{Not} RFC 6979 / $k_\text{proto}$ \\
         SSH Term & 7.0.26 & \faApple & \unsupported & \unsupported & \unsupported & \unsupported & OpenSSL 3.5.0 & \textit{Not} RFC 6979 / $k_\text{proto}$ \\
         Tectia SSH & 6.6.3.490 & \faWindows & \unsupported & \supported & \supported & \faBolt & \textit{Proprietary} & Random \\
         ~ & ~ & ~ & ~ & ~ & ~ & ~ & OpenSSL 3.0.8 (FIPS)& Random \\
         Tera Term & 5.2 & \faWindows & \supported & \unsupported & \supported & \supported & LibreSSL 3.8.2 & Random \\
         Termius & 9.8.5 & \faLinux & \unsupported & \supported & \supported & \supported & Libsodium 1.0.17 & Random \\
         ~ & 9.21.2 & \faApple & \unsupported & \supported & \supported & \supported & Libsodium 1.0.17 & Random \\
         Win32 OpenSSH & 9.5.0.0 & \faWindows & \supported & \supported & \supported & \supported & LibreSSL 3.8.2 & Random \\
         WinSCP & 6.3.4 & \faWindows & \supported & \unsupported & \supported & \supported & \putty \puttyfixedversion & RFC 6979 \\
         XShell 7 & 0170 & \faWindows & \unsupported & \supported & \supported & \supported & OpenSSL 1.0.2u & Random \\
         \bottomrule
    \end{tabular}
\end{table*}

\begin{table}
    \centering
    \caption{Overview of SSH agent implementations, their underlying cryptographic library, and nonce generation method. 1Password does not support any (EC)DSA algorithms.}
    \label{tab:agents}
    \begin{tabular}{llll}
        \toprule
         Agent & Version & Crypto Library & Nonce  \\ 
         \midrule
         1Password & 8.10.56 & Libgcrypt 1.10.3 & \textit{n/a} \\
         GnuPG & 2.4.4 & Libgcrypt 1.10.3 & RFC 6979 \\
         Goldwarden & 0.3.6 & Go 1.22.1 crypto & Random \\
         KeeAgent & 0.13.8 & BC C\# 1.9.0 & Random \\
         MobaXTerm & 24.4 & \putty \puttyfixedversion & RFC 6979 \\
         OpenSSH & 9.6p1 & OpenSSL 3.0.13 & Random \\
         PKIX-SSH & 15.3 & OpenSSL 3.0.13 & Random \\
         \putty Pageant & \puttyvulnversion & \putty \puttyvulnversion & $k_\text{proto}$ \\ 
         ~ & \puttyfixedversion & \putty \puttyfixedversion & RFC 6979 \\
         SecureCRT & 9.5.2 & \makecell[l]{\scalebox{.67}[1.0]{BSAFE Crypto-C 
         ME 4.1.5}} & Random \\
         Tectia SSH & 6.6.3.490 & \textit{Proprietary} & Random \\
         ~ & ~ & OpenSSL 3.0.8 & Random \\
         Termius & 9.9.0 & Libsodium 1.0.17 & Random \\
         Win32 OpenSSH & 9.5.0.0 & LibreSSL 3.8.2 & Random \\
         XShell 7 Xagent & 0170 & OpenSSL 1.0.2u & Random \\
         \bottomrule
    \end{tabular}
\end{table}

\subsection{\putty P-521 Private Key Recovery}
\label{subsec:putty}

\begin{figure}
	\centering
    \includegraphics[width=\columnwidth]{img/SuccessRateBiasedNoncePuTTY.pdf}
    \caption[Success rate for the biased nonce attack on \putty ECDSA P-521.]{
      Success rate of the biased nonce attack on \putty ECDSA \pfivetwoone using USVPPredSieve~\cite{EC:AlbHen21}, for a given number of signatures (1024 trials each).
    }
    \Description{A line chart depicting the success rate of biased nonce attack on \putty ECDSA P-521 in relation to the available signature count. At 58 signatures, the success rate is greater than 0.5. At 60 signatures, the success rate is 1.}
    \label{fig:puttyeval}
\end{figure}

\paragraph{\putty's Deterministic Nonce Generation ($k_\text{proto}$)} 

\putty generates nonces deterministically to protect against bad random number generation that could result in biased nonces. The initial implementation in~2001 was designed for DSA, but not ECDSA, and generated the nonce by reducing a 512-bit pseudorandom number $k_\text{proto}$ to a number in the interval $[2,q-1]$. $k_\text{proto}$ is derived from a label, the private key $x$, and the message digest $\text{H}(m)$.\footnote{\url{https://github.com/github/putty/blob/7003b43963aef6cdf2841c2a882a684025f1d806/sshdss.c\#L337}, accessed: 2025-04-14}

\begin{align*}
  \text{digest} &:= \text{SHA-512}\,(\mathit{label} \,||\, x) \\
  k_\text{proto} & := \text{SHA-512}\left(\text{digest} \,||\, \text{H}(m)\right) \\
  k &:= \left(k_\text{proto} \bmod q-2\right) + 2
\end{align*}

For DSA, $q$ is the order of the cyclic subgroup of $\mathbb{Z}_p^*$ generated by the basis element $g$ and is 160 bits in the case of SSH. When support for ECDSA was added in~2014, the code for $k_\text{proto}$ was reused, except for using a different label. In ECDSA, $q$ is the order of the elliptic curve used for all computations generated by the base point $P$.

\paragraph{Bias in P-521}

For DSA, and ECDSA with NIST~\ptwofivesix and \pthreeeightfour, the bias on the nonces $k$ generated with $k_\text{proto}$ is negligible. However, for \pfivetwoone, the output of \shafiveonetwo is always smaller than the size of an element of the underlying field $GF(2^{521}-1)$, and thus the value $k_\text{proto}$ is not reduced. As a result, the nonce $k < 2^{512}$ has the upper nine most significant bits set to zero. 

\paragraph{Hidden Number Problem}

In ECDSA, this nonce $k$ is used in the generation of the signature $(r,s)$ on a message $m$ using the private key $a$ as follows: $r = \text{pr}_x(kP)$ and $s = k^{-1} \cdot (h(m)+a \cdot r) \bmod q$, where $\text{pr}_x$ is the x-coordinate of a point. Rewriting this yields
$$
s^{-1} \cdot r \cdot a + s^{-1}\cdot h(m) \equiv k   \pmod q.
$$
Setting $t = s^{-1} \cdot r$, $b = s^{-1}\cdot h(m)$,  and $\alpha = a$ results in the equivalence relation of the \emph{Hidden Number Problem} (HNP), as formulated by Boneh and Venkatesan~\cite{C:BonVen96},
$$
t \cdot \alpha + b \equiv k \pmod q.
$$
Here, $b$ and $t$ are values known to the attacker, the hidden number $\alpha$ is the secret key the attacker wants to recover, and $k$ is partially known to the attacker, as its nine most significant bits are $0$. Given ``enough'' equivalences, the HNP can be solved, i.e., the value $\alpha = a$ can be computed. The number of HNP equivalences needed to solve the HNP problem (i.e., to compute the private ECDSA signing key) depends on the algorithm used.
 
\paragraph{Lattice Algorithms to Solve the HNP}

For our experiments, we used an algorithm based on lattices from~\cite[USVPPredSieve]{EC:AlbHen21} (see also~\cite{FC:BreHen19}), which can be directly applied to our setting and recovers the private key by solving the Shortest Vector Problem in the constructed lattice. The results are shown in~\autoref{fig:puttyeval}. Given 58 ECDSA signatures generated with $k_\text{proto}$ and \pfivetwoone, we can compute the signer's private key with a probability greater than 0.5. Given 60 such signatures, we can \emph{always} compute the private signing key. We repeated the experiment 1024 times for each signature count, and the computation time for each trial was at most 67 seconds.

\paragraph{Root Cause Analysis}

For the curves \poneninetwo, \ptwotwofour, \ptwofivesix, and \pthreeeightfour, the bit lengths of the coordinates of the curve points were carefully aligned both with byte boundaries and with the output lengths of hash functions: Both hash function families \shatwo and \shathree have hash value lengths of 224, 256, and 384 bits, and the output length 192 appears in NIST SP 800\nobreakdashes-208 as a truncated version of \shatwofivesix. The only exception to this rule is \pfivetwoone, for using a Mersenne prime. Informed programmers wonder about possible typos (512 vs. 521), and discuss why this parameter was chosen.\footnote{\url{https://crypto.stackexchange.com/questions/6219/why-do-the-elliptic-curves-recommended-by-nist-use-521-bits-rather-than-512}, accessed: 2025-04-14}
